\chapter{Internes}
\label{chp:INT}

\section{Ausbilungsverordnung}

Ausbildungsziele für österreichische Rettungssanitäter gem. der \emph{Verordnung
der Bundesministerin für Gesundheit und Frauen über die Ausbildung zum
Sanitäter - Sanitäter-Ausbildungsverordnung - San-AV StF: BGBl. II Nr. 420/2003
-- Anlage 1, Modul I}.

% \clearpage
%%%%%%%%%%%%%%%
\tablefirsthead{%
\toprule
\multicolumn{1}{c}{\bfseries Inhalt} &

\multicolumn{1}{c}{\bfseries Referenz} 
\\
\midrule}
%%%%%%%%%%%%%%%
\tablehead{%
\toprule
\multicolumn{2}{l}{\small\itshape Fortsetzung.}
\\
\midrule
\multicolumn{1}{c}{\bfseries Inhalt} &
\multicolumn{1}{c}{\bfseries Referenz} 
\\
\midrule}
%%%%%%%%%%%%%%%
\tabletail{%
\hline
\multicolumn{2}{r}{\small\itshape weiter \ldots}
\\
\bottomrule}
%%%%%%%%%%%%%%%
\tablelasttail{\bottomrule}
\bottomcaption{This table is split across pages}

\begin{xtabular}{p{10cm}p{5cm}}

\gTabHeading{Erste Hilfe und erweiterte Erste Hilfe}  

(Bestandteil des Sachgebiets 'Sanitätshilfe') 

\begin{inparadesc}
\item[Stundenanzahl:] 14
\item[Lehrkraft:] Notarzt, Arzt für Allgemeinmedizin, approbierter Arzt,
Facharzt, Lehrsanitäter, fachkompetente Person  
\end{inparadesc}
\\
- Notwendigkeit, Verpflichtung der Hilfeleistung
& 
\ref{}\noindent\gCa{k.A.}
\\
- Unfallverhütung          
& 
\noindent\gCa{k.A.}
\\
- Rettungskette           
& 
\ref{}\noindent\ref{topic:rettungskette}
\\
- Gefahrenzonen            
& 
\noindent\ref{topic:gefahrenzone}
\\
 - Bergen (Wegziehen, Rautekgriff aus dem  Auto, Sturzhelm)         
& 
\noindent\ref{topic:wegziehen}\ref{topic:rautekgriff}\ref{topic:sturzhelmabnahme}
\\
 - Notfallpatient, Notfalldiagnose, Notfallhilfe
& 
\noindent\ref{topic:notfallpatient}\ref{topic:notfalldiagnose}\ref{topic:notfallhilfe}
\\
 - Kontrolle der Lebensfunktionen
& 
\noindent\ref{topic:lebensfunktionen-kontrolle}
\\
 - Notfalldiagnose Bewusstlosigkeit
& 
\ref{topic:notfalldiagnose-bewusstlosigkeit}
\\
- Notfalldiagnose  Atem-Kreislaufstillstand
& 
\noindent\ref{topic:notfalldiagnose-atem-kreislaufstillstand}
\\
- Starke Blutung (Blutstillung, Fingerdruck,  Druckverband, Abbindung)
& 
\noindent\ref{topic:}
\\
 - Schock (Schockbekämpfung,  Lagerungen)
& 
\noindent\gCa{k.A.}
\\
- Wunden und Wundverbände (Tierbisse, Insektenstiche, Verätzungen,Verbrennungen,Erfrierungen,Unterkühlung)
& 
\noindent\gCa{k.A.}
\\
- Quetschungen
& 
\noindent\gCa{k.A.}
\\
- Gelenksverletzungen
& 
\noindent\gCa{k.A.}
\\
- Knochenbrüche, allgemein
& 
\noindent\ref{topic:knochenbruch}
\\
- Brustkorbverletzungen
& 
\noindent\gCa{k.A.}
\\
- Plötzlich auftretende Erkrankungen
& 
\noindent\gCa{k.A.}
\\
- Vergiftungen
& 
\noindent\gCa{k.A.}
\\
- Transport
& 
\noindent\gCa{k.A.}
\\\toprule
\gTabHeading{Hygiene}  

(Bestandteil des Sachgebiets 'Sanitätshilfe') 

\begin{inparadesc}
\item[Stundenanzahl:] 2
\item[Lehrkraft:] Arzt für Allgemeinmedizin, approbierter Arzt,
Facharzt, Angehörige des gehobenen Dienstes für Gesundheits- und Krankenpflege,
Hygienefachkraft,  Lehrsanitäter
\end{inparadesc}

\\\midrule
- Persönliche Hygiene      
& 
\ref{topic:persoenliche-hygiene}
\\
- Grundbegriffe der Desinfektion             
& 
\ref{topic:desinfektion-grundzuege}
\\
- Grundbegriffe der  Sterilisation         
& 
\ref{topic:sterilisatiom-grundzuege}
\\
- Entsorgung von infektiösem Abfall       
& 
\ref{}\noindent\gCa{k.A.}
\\
- Allgemeine Infektionslehre          
& 
Erreger: \ref{topic:infektion-erreger}, Übertragungswege:
\ref{topic:infektion-uebertragungswege}, Kreuzinfektion:
\ref{topic:kreuzinfektion}.
\\
- Vorgehen bei Verletzungen des Personals                
& 
Nadelstichverletzungen: \ref{topic:nadelstichverletzungen}
\\
- Hygienemaßnahmenplan     
& 
\ref{topic:hygieneplan}
\\
- Infektionstransport      
& 
\gAbk{MRSA} und andere resistente Keime: \ref{topic:mrsa-transport},
sonstige: \noindent\gCa{k.A.} \\\toprule
\gTabHeading{Berufsspezifische rechtliche Grundlagen}  ~



\begin{inparadesc}
\item[Stundenanzahl:] 3
\item[Lehrkraft:]  Jurist, fachkompetente Person
\end{inparadesc}

\\\midrule

- Aufgaben und Kompetenzen des  Rettungssanitäters       
& 
\ref{}\noindent\gCa{k.A.}
\\
- Dokumentation im Rettungswesen (Einsatzprotokoll, Leitstellendokumentation, Transportnachweis)
& 
\ref{}\noindent\gCa{k.A.}
\\
- Hilfs- und Rettungswesen
& 
\ref{}\noindent\gCa{k.A.}
\\
- Straßenverkehrsordnung
& 
\noindent\ref{topic:strassenverkehrsordnung}
\\
- Patientenrechte
& 
\noindent\ref{topic:patientenrechte}
\\
- Grundlagen des Haftungsrechtes
& 
\noindent\ref{topic:haftungsrecht}
\\
- Unterbringungsgesetz
& 
\noindent\ref{topic:unterbringungsgesetz}
\\
- Reversfähigkeiten und Effekten
& 
\noindent\ref{topic:reversfähigkeit}\ref{topic:effekte}
\\
- Mitnahme von Begleitpersonen
& 
\noindent\ref{topic:begleitpersonen-mitnahme}
\\\toprule
\gTabHeading{Anatomie und Physiologie}  

(Bestandteil des Sachgebiets 'Sanitätshilfe')  

\begin{inparadesc}
\item[Stundenanzahl:] 4
\item[Lehrkraft:]  Notarzt, Arzt für Allgemeinmedizin, approbierter Arzt,
Facharzt, Turnusarzt,   Angehörige  des  gehobenen  Dienstes für Gesundheits- 
und Krankenpflege, Lehrsanitäter, fachkompetente Person
\end{inparadesc}

\\\midrule

- Blutkreislauf - Grundzüge               
& 
\ref{topic:blutkreislauf}
\\
- Gliedmaßen - Grundzüge  
& 
Obere Extremität: \ref{topic:schulterguertel}, \ref{topic:obere-extremitaet},
untere Extremität: \ref{topic:beckenguertel}, 
\ref{topic:untere-extremitaet}
\\
- Haut - Grundzüge        
& 
\ref{topic:haut}
\\
- Schädel und Rumpf - Grundzüge Skelett       
& 
Schädel: \ref{topic:schaedel}
\\
- Brustkorb - Grundzüge   
& 
\ref{topic:thorax}
\\
- Bauchraum - Grundzüge   
& 
Gastro-intestinaltrakt: \ref{topic:gastrointestinaltrakt}, andere Bauchorgane:
\ref{topic:sonstige-bauchorgane}. 
\\\toprule
\gTabHeading{Störungen der  Vitalfunktionen  und Regelkreise und zu setzende Maßnahmen}  

(Bestandteil des  Sachgebiets  'Sanitätshilfe') 

\begin{inparadesc}
\item[Stundenanzahl:]  8 
\item[Lehrkraft:]  Notarzt, Arzt für Allgemein- medizin, approbierter Arzt,
Facharzt,  Lehr- sanitäter,  fachkompe- tente Person
\end{inparadesc}

\\\midrule

- Definition Vitalfunktion           
& 
\noindent\ref{topic:vitalfunktionen-definition}
\\
- NACA-Schema             
& 
\ref{topic:naca}
\\
- Bewusstsein             

(Bewusstseinsstörungen, 
& 
\noindent Bewusstsein
allgemein: \ref{topic:bewusstsein};
Bewusstseinsstörungen: \ref{topic:bewusstseinsstoerung}
\\
Bewusstlosigkeit)       
& 
\noindent\ref{topic:bewusstlosigkeit}
\\
- Atmung 
& 
--- ↓ --- 
\\          
(Atemstörungen,  
& 
\noindent\ref{topic:atemstörungen}
\\
Atemstillstand)        
& 
\noindent\ref{topic:atemstillstand}
\\
- Kreislauf               
& 
--- ↓ --- 
\\
(Kreislaufstörungen,    
& 
\ref{}\noindent\gCa{k.A.}
\\
Kreislaufstillstand)    
& 
\ref{}\noindent\gCa{k.A.}
\\
- Regelkreise
& 
--- ↓ --- 
\\
(Wärmehaushalt, 
& 
\ref{}\noindent\gCa{k.A.}
\\
Wasser- und Elektrolythaushalt,
& 
\ref{topic:wasser-elektrolythaushalt}
\\
Säure-Basen-Haushalt,
& 
\ref{topic:saeure-basen-haushalt}
\\
Stoffwechsel)
& 
\ref{topic:stoffwechsel}
\\
- Starke Blutung
& 
\ref{}\noindent\gCa{k.A.}
\\
- Schock (Ursachen,
& 
--- ↓ --- \ref{topic:schock}
\\
Wirkung, Schockformen,
& 
\ref{topic:schock-arten}
\\
Verlauf, Schockzeichen)
& 
\ref{}\noindent\gCa{k.A.}
\\
- Akute Störung der Atmung 
& 
--- ↓ --- 
\\
(Atembehinderung, 
& 
\ref{}\noindent\gCa{k.A.}
\\
Verlegung der Atemwege  durch Fremdkörper, 
& 
\ref{}\noindent\gCa{k.A.}
\\
Verlegung durch Schwellung)
& 
\ref{}\noindent\gCa{k.A.}
\\
- Akute Atemnot
& 
\ref{}\noindent\gCa{k.A.}
\\
- Absaugung
& 
\ref{}\noindent\gCa{k.A.}
\\
- Sauerstoff (Inhalationsrichtlinien, Dosierungsvorschriften)
& 
\ref{topic:sauerstoff}
\\
- Assistierte Beatmung (Indikationen, Durchführung)
& 
\ref{}\noindent\gCa{k.A.}
\\
- Feststellung des Todes
& 
\ref{topic:tod}
\\\toprule
\gTabHeading{Notfälle bei verschiedenen Krankheitsbildern  und zu setzende Maßnahmen}  

(Bestandteil des Sachgebiets 'Sanitätshilfe') 

\begin{inparadesc}
\item[Stundenanzahl:]  6
\item[Lehrkraft:]  Notarzt, Arzt für Allgemeinmedizin, approbierter Arzt,
Facharzt, Turnusarzt, Lehrsanitäter, fachkompetente Person
\end{inparadesc}
\\\midrule
- Koma aus vorerst unbekannter Ursache      
& 
--- ↓ --- 
\\
(Schlaganfall,           
& 
\ref{topic:insult}
\\
Meningitis, 
& 
\ref{topic:meningitis}
\\
Diabetes,   
& 
\ref{}\noindent\gCa{k.A.}
\\
Vergiftung)             
& 
\ref{}\noindent\gCa{k.A.}
\\
- Krampfanfall             
& 
--- ↓ --- 
\\
(Epilepsie, Tetanie)    
& 
\ref{}\noindent\gCa{k.A.}
\\
- Akuter Gefäßverschluss an den Gliedmaßen     
& 
\ref{}\noindent\gCa{k.A.}
\\
(Venenthrombose,       
& 
\ref{}\noindent\gCa{k.A.}
\\
Arterielle Embolie)     
& 
\ref{}\noindent\gCa{k.A.}
\\
- Pulmonale Notfälle      
& 
--- ↓ ---
\\
(Asthma bronchiale,
& 
\ref{}\noindent\gCa{k.A.}
\\
Lungenödem, 
& 
\ref{}\noindent\gCa{k.A.}
\\
Lungenembolie, 
& 
\ref{}\noindent\gCa{k.A.}
\\
Lungenentzündung)
& 
\ref{}\noindent\gCa{k.A.}
\\
- Cardiale Notfälle 
& 
--- ↓ --- 
\\
(Herz-Rhythmus-Störungen, 
& 
\ref{}\noindent\gCa{k.A.}
\\
Herzversagen, Linksherzschwäche,
Rechtsherzschwäche, 
& 
\ref{}\noindent\gCa{k.A.}
\\
Angina pectoris, 
& 
\ref{}\noindent\gCa{k.A.}
\\
Herzinfarkt, 
& 
\ref{}\noindent\gCa{k.A.}
\\
Hochdruckkrise) 
& 
\ref{}\noindent\gCa{k.A.}
\\
- Allgemeinchirurgische Notfälle & --- ↓ --- 
\\
(Akutes Abdomen, & \ref{topic:akutes-abdomen}
\\
Gastritis, & \noindent\gCa{k.A.}
\\
Ulcus, & \noindent\gCa{k.A.}
\\
Gastroenteritis,
& 
\noindent\gCa{k.A.}
\\
Pankreatitis,
& 
\noindent\gCa{k.A.}
\\
Appendizitis,
& 
\ref{topic:appendizitis}
\\
Gallenkolik, 
& 
\ref{topic:gallenkolik}
\\
Ileus,
& 
\noindent\gCa{k.A.}
\\
Mesenterialinfarkt,
& 
\noindent\gCa{k.A.}
\\
Lebensmittelvergiftung,
& 
\noindent\gCa{k.A.}
\\
gastrointestinale Blutung)
& 
\noindent\gCa{k.A.}
\\
- Gynäkologische und urologische Erkrankungen
& 
--- ↓ --- 
\\
Nierenkolik,
& 
\ref{topic:nierenkolik}
\\
Nierenbeckenentzündung,
& 
\noindent\gCa{k.A.}
\\
Harnwegsinfekt, 
& 
\noindent\gCa{k.A.}
\\
akutes Harnverhalten,
& 
\noindent\gCa{k.A.}
\\
chronische Niereninsuffienz und Hämodialyse,
& 
\noindent\gCa{k.A.}
\\
extrauterine Gravidität, 
& 
\noindent\gCa{k.A.}
\\
Abortus,
& 
\noindent\gCa{k.A.}
\\
Unterleibsblutungen,
& 
\noindent\gCa{k.A.}
\\
Ovarialtumore,
& 
\noindent\gCa{k.A.}
\\
Vergewaltigung (typische Verletzungen, psychische Betreuung)
& 
\noindent\gCa{k.A.}
\\\toprule
\gTabHeading{Spezielle Notfälle und zu setzende Maßnahmen  }  

(Bestandteil des Sachgebiets  'Sanitätshilfe')  

\begin{inparadesc}
\item[Stundenanzahl:]  15  
\item[Lehrkraft:]   Notarzt, Arzt für  Allgemein-  medizin,  approbierter 
Arzt, Facharzt,   Turnusarzt,  Lehr-  sanitäter,  fachkompetente Person
\end{inparadesc}

\\\midrule

- Traumatologische Notfälle
& 
--- ↓ --- 
\\ 
(Schädel-Hirn- Verletzungen,          
& 
\noindent\gCa{k.A.}
\\
Hirnblutung 
& 
\noindent\gCa{k.A.}
\\
Hirndruck) 
& 
\noindent\gCa{k.A.}
\\
- Halswirbelsäulen- und  Wirbelsäulentrauma     
& 
--- ↓ --- 
\\
(Sturzhelmabnahme,      
& 
\noindent\gCa{k.A.}
\\
Halswirbelsäulenschienung, 
& 
\noindent\gCa{k.A.}
\\
Motorik-Durchblutung-Sensibilitätskontrolle,         
& 
\noindent\gCa{k.A.}
\\
Body-Check, 

Umgang mit  Schaufeltrage und Vakuummatratze,

Sandwich-Technik)

- Extremitätentrauma

--- ↓ --- 

(Verletzungsarten,

Motorik-Durchblutung-Sensibilitätskontrolle,

Stiefelgriff, Prinzip der Schienung,

Ruhigstellung,

Schienung des Armes,
& 
\ref{}\noindent\gCa{k.A.}
\\
Schienung des Beines,
& 
\ref{}\noindent\gCa{k.A.}
\\
Pneumatische Schiene,
& 
\ref{}\noindent\gCa{k.A.}
\\
Vakuumschiene,
& 
\ref{}\noindent\gCa{k.A.}
\\
Extensionsschiene)
& 
\ref{}\noindent\gCa{k.A.}
\\
- Thoraxtrauma 
& 
--- ↓ --- \ref{topic:thoraxtrauma}
\\
(offene, geschlossene Brustkorbverletzung,
& 
\ref{}\noindent\gCa{k.A.}
\\
Serienrippenbruch,
& 
\ref{topic:serienrippenfraktur}
\\
Pneumothorax,
& 
\ref{topic:pneumothorax}
\\
Spannungspneumothorax)
& 
\ref{topic:pneumothorax}
\\
- Bauchtrauma 
& 
\ref{topic:bauchtrauma}
\\
(offene, geschlossene Bauchverletzung,
& 
\ref{}\noindent\gCa{k.A.}
\\
Verletzung der Harn und Geschlechtsorgane)
& 
\ref{}\noindent\gCa{k.A.}
\\
- Beckentrauma
& 
\ref{topic:beckentrauma}
\\
- Polytrauma (Definition, Prioritäten, Management)
& 
\ref{topic:polytrauma}
\\
- Wunden (mechanische, chemische, thermische)
& 
\ref{topic:wunden}
\\
- Dekubitus Prophylaxe und Lagerung bei Dekubitus
& 
\ref{topic:}\noindent\gCa{k.A.}
\\
- Verbandlehre
& 
\ref{topic:}\noindent\gCa{k.A.}
\\
- Akut auftretende Blutungen
& 
\ref{topic:}\noindent\gCa{k.A.}
\\
- Vergiftung (Ursachen und Verdacht, Aufnahmearten)
& 
\ref{topic:}\noindent\gCa{k.A.}
\\
- Psychiatrische Notfälle
& 
--- ↓ --- 
\\
(Suizid, 
& 
\ref{topic:suizid}
\\
Psychose,
& 
\ref{topic:}\noindent\gCa{k.A.}
\\
Suchterkrankungen und Entzugssyndrom,
& 
\ref{topic:}\noindent\gCa{k.A.}
\\
Depression, 
& 
\ref{topic:depression}
\\
Manie)
& 
\ref{topic:manie}
\\
- Schwangerschaft und Geburt 

(Notfälle in der Schwangerschaft,
& 
\ref{topic:}\noindent\gCa{k.A.}
\\
Geburt,
& 
\ref{topic:geburt}
\\
Geburtskomplikationen,
& 
\ref{topic:}\noindent\gCa{k.A.}
\\
Versorgung des Neugeborenen)

- Notfälle im Säuglings- und Kleinkindalter
& 
--- ↓ --- 
\\
(anatomische und physiologische Besonderheiten,
& 
\ref{topic:}\noindent\gCa{k.A.}
\\
Kontrolle der Lebensfunktionen und lebensrettende Sofortmaßnahmen,
& 
\ref{topic:}\noindent\gCa{k.A.}
\\
Krampfanfälle,
& 
Allgemein: \ref{topic:krampfanfaelle}, Fieberkrampf: \label{topic:fieberkrampf}
\\
Pseudokrupp,
& 
\ref{topic:}\noindent\gCa{k.A.}
\\
Epiglottitis,
& 
\ref{topic:epiglottitis}
\\
Keuchhusten,
& 
\ref{topic:}\noindent\gCa{k.A.}
\\
plötzlicher Kindstod)
& 
\ref{topic:sids}
 \\\toprule
\gTabHeading{Defibrillation mit halbautomatischen Geräten}  

(Bestandteil des Sachgebiets  'Sanitätshilfe')  

\begin{inparadesc}
\item[Stundenanzahl:]   8
\item[Lehrkraft:]   Notarzt,  Arzt für Allgemeinmedizin,  approbierter Arzt,
Facharzt, Lehrsanitäter, fachkompetente Person
\end{inparadesc}

\\\midrule
- Der halbautomatische Defibrillator          

- Handhabung eines halbautomatischen Defibrillators         

- Gerätemanagement während der Reanimation 

- Erfolgskontrolle
 \\\toprule
\gTabHeading{Gerätelehre und Sanitätstechnik }  



\begin{inparadesc}
\item[Stundenanzahl:]  12 
\item[Lehrkraft:]   Lehrsanitäter, fachkompetente Person
\end{inparadesc}

\\\midrule

- Medizinproduktegesetz - Information    
& 
\ref{topic:medizinproduktegesetz}
\\
- Bergungs- und Lagerungstechniken     
& 
\ref{}\noindent\gCa{k.A.}
\\
(Bergetuch,
& 
\ref{topic:bergetuch}
\\
Einheitskrankentrage,
& 
\ref{}\noindent\gCa{k.A.}
\\
Tragsessel, 
& 
\ref{topic:tragsessel}
\\
Fahrtrage,
& 
\ref{topic:fahrtrage}
\\
Rollstuhl)
& 
\ref{}\noindent\gCa{k.A.}
\\
- Einsatzfahrzeug
& 
\ref{}\noindent\gCa{k.A.}
\\
- Beatmungsbeutel
& 
\ref{topic:beatmungsbeutel}
\\
- Absauggeräte
& 
\ref{topic:absaugung}
\\
- Sauerstoff
& 
\ref{topic:sauerstoff}
\\
- Infusionen und
& 
\ref{}\noindent\gCa{k.A.}
\\
Infusionsgeräte
& 
\ref{}\noindent\gCa{k.A.}
\\
- Blutdruckmessung
& 
\ref{topic:rr-messung}
\\
- Stabilisierungs- und Schienungstechniken
& 
\ref{}\noindent\gCa{k.A.}
\\
(Stabilisierung der Halswirbelsäule,
& 
\ref{}\noindent\gCa{k.A.}
\\
Schaufeltrage,
& 
Allg.: \ref{topic:schaufeltrage}, Anwendung:
\ref{topic:schaufeltrage-anwendung}
\\
Vakuummatratze)
& 
Allg.: \ref{topic:vakuummatratze}, Anwendung:
\ref{topic:vakuummatratze-anwendung}
\\
- Geburtenkoffer
& 
\ref{}\noindent\gCa{k.A.}
\\
- Transportinkubator
& 
\ref{}\noindent\gCa{k.A.}
\\
 \\\toprule
\gTabHeading{Rettungswesen}  



\begin{inparadesc}
\item[Stundenanzahl:]  4  
\item[Lehrkraft:]   Fachkompetente Person
\end{inparadesc}

\\\midrule

- Rechtliche Grundlagen  
& 
\ref{}\noindent\gCa{k.A.}
\\
- Zusammenarbeit mit  anderen Organisationen
& 
\ref{}\noindent\gCa{k.A.}
\\
- Einsatzarten
& 
\ref{}\noindent\gCa{k.A.}
\\
- Transport- und Fahrzeugarten (Land,  Wasser, Luft)
& 
\ref{}\noindent\gCa{k.A.}
\\
- Normen, persönliche Schutzausrüstung
& 
\ref{}\noindent\gCa{k.A.}
\\
- Fahrzeugausstattung
& 
\ref{}\noindent\gCa{k.A.}
\\
- Rettungskette, Hilfsfrist
& 
\ref{}\noindent\gCa{k.A.}
\\
- Dienststellennetz
& 
\ref{}\noindent\gCa{k.A.}
\\
- Personal im Rettungsdienst
& 
\ref{}\noindent\gCa{k.A.}
\\
- Notarztsysteme
& 
\ref{}\noindent\gCa{k.A.}
\\
- Leitstelle, Kommunikationsarten
& 
\ref{}\noindent\gCa{k.A.}
\\
- Gefahren an der Einsatzstelle
& 
\ref{}\noindent\gCa{k.A.}
\\
- Gefahrguteinsätze, Sondertransporte

 \\\toprule
\gTabHeading{Katastrophen, Großschadensereignisse, Gefahrgutunfälle}  



\begin{inparadesc}
\item[Stundenanzahl:]  4  
\item[Lehrkraft:]  Fachkompe- tente Person
\end{inparadesc}

\\\midrule

- Katastrophen  
& 
\ref{}\noindent\gCa{k.A.}
\\
(Rechtliche Grundlagen, 
& 
\ref{}\noindent\gCa{k.A.}
\\
 Geltungsbereiche, Arten der Katastrophen,
& 
\ref{}\noindent\gCa{k.A.}
\\
Phasen der Katastrophenbewältigung,
& 
\ref{}\noindent\gCa{k.A.}
\\
Katastrophenhilfseinheiten,
& 
\ref{}\noindent\gCa{k.A.}
\\
Führungsorganisation,
& 
\ref{}\noindent\gCa{k.A.}
\\
personelle, materielle und finanzielle Vorsorge,
& 
\ref{}\noindent\gCa{k.A.}
\\
Einsatzgrundsätze,
& 
\ref{}\noindent\gCa{k.A.}
\\
generelle Einsatzrichtlinien,
& 
\ref{}\noindent\gCa{k.A.}
\\
Grundzüge der Triage)
& 
\ref{topic:triage}
\\
- Großschadensereignisse
& 
\ref{}\noindent\gCa{k.A.}
\\
(Rechtliche Grundlagen,
& 
\ref{}\noindent\gCa{k.A.}
\\
Einstufung,
& 
\ref{}\noindent\gCa{k.A.}
\\
Alarmierung,
& 
\ref{}\noindent\gCa{k.A.}
\\
Schadensraum,
& 
\ref{}\noindent\gCa{k.A.}
\\
Schadensplatz, Sicherheitseinrichtungen,
& 
\ref{}\noindent\gCa{k.A.}
\\
Organisation beim Großunfall, 
& 
\ref{}\noindent\gCa{k.A.}
\\
Grundzüge der Triage,
& 
\ref{topic:triage}
\\
österreichisches Patientenleitsystem,
& 
\ref{topic:pls}
\\
Material und Ausrüstung,
& 
\ref{}\noindent\gCa{k.A.}
\\
Kommunikation)
& 
\ref{}\noindent\gCa{k.A.}
\\
- Gefahrgutunfälle 
& 
\ref{}\noindent\gCa{k.A.}
\\
(Arten von Gefahrgutunfällen,
& 
\ref{}\noindent\gCa{k.A.}
\\
Gefahrzettel,
& 
\ref{topic:gefahrzettel}
\\
Gefahrsymbole,
& 
\ref{topic:gefahrsymbole}
\\
Warntafel, 
& 
\ref{topic:warntafel}
\\
Verhalten am Unfallort, 
& 
\ref{}\noindent\gCa{k.A.}
\\
Koordination mit anderen Einsatzorganisationen,
& 
\ref{}\noindent\gCa{k.A.}
\\
Absperrmaßnahmen,
& 
\ref{}\noindent\gCa{k.A.}
\\
Sofortmaßnahmen)
& 
\ref{}\noindent\gCa{k.A.}
\\
 \\\toprule
\gTabHeading{Angewandte Psychologie und Stressbewältigung}  



\begin{inparadesc}
\item[Stundenanzahl:]  4 
\item[Lehrkraft:]  Fachkompe- tente Person
\end{inparadesc}

\\\midrule

- Belastung, Anforderung, Beanspruchung, workflow
& 
\ref{}\noindent\gCa{k.A.}
\\
- Überforderung, Unterforderung
& 
\ref{}\noindent\gCa{k.A.}
\\
- Beanspruchungsfolgen
& 
\ref{}\noindent\gCa{k.A.}
\\
- Stressursachen, -entstehung und -faktoren
& 
\ref{}\noindent\gCa{k.A.}
\\
- Stressauswirkung
& 
\ref{}\noindent\gCa{k.A.}
\\
- Früherkennung
& 
\ref{}\noindent\gCa{k.A.}
\\
- Grundsätze der Stressvermeidung
& 
\ref{}\noindent\gCa{k.A.}
\\
- Maßnahmen zur Verhütung und Verminderung von Beanspruchungsfolgen
& 
\ref{}\noindent\gCa{k.A.}
\\
- Psychische Betreuung von Kranken/Verletzten
& 
\ref{}\noindent\gCa{k.A.}
\\
(Gesprächsführung, 
& 
\ref{}\noindent\gCa{k.A.}
\\
Vertrauensaufbau und Patienteninformation, 
& 
\ref{}\noindent\gCa{k.A.}
\\
psychische Belastungssyndrome, 
& 
\ref{}\noindent\gCa{k.A.}
\\
verwirrte Patienten, 
& 
\ref{}\noindent\gCa{k.A.}
\\
Begleitung und Betreuung Sterbender, 
& 
\ref{}\noindent\gCa{k.A.}
\\
Supervision)
& 
\ref{}\noindent\gCa{k.A.}
\\
 \\\toprule
\gTabHeading{Praktische Übungen ohne Patientenkontakt}  



\begin{inparadesc}
\item[Stundenanzahl:] 16
\item[Lehrkraft:] Lehrsanitäter, fachkompetente Person
\end{inparadesc}

\\\midrule

- Regloser Notfallpatient  
& 
\gTakeNotes
\\
- Kontrolle der  Lebensfunktionen  (erwachsener Notfallpatient)
& 
\gTakeNotes
\\
- Notfalldiagnose Bewusstlosigkeit (stabile Seitenlage)
& 
\gTakeNotes
\\
- Notfalldiagnose Atemstillstand (Beatmung)
& 
\gTakeNotes
\\
- Notfalldiagnose Kreislaufstillstand (Beatmung und Herzmassage)
& 
\gTakeNotes
\\
- Blutstillung (Fingerdruck, Abdrückstellen, Druckverband, Abbindung, Amputatsversorgung)
& 
\gTakeNotes
\\
- Blutdruckmessung
& 
\gTakeNotes
\\
- Schockbekämpfung (Lagerungsarten)
& 
\gTakeNotes
\\
- Halswirbelsäulen- und Wirbelsäulentrauma 
&
--- ↓ --- 
\\
(Sturzhelmabnahme, 
& 
\gTakeNotes
\\
Halswirbelsäulenschienung, 
& 
\gTakeNotes
\\
Motorik-Durchblutung-Sensibilitätskontrolle, 
& 
\gTakeNotes
\\
Body-Check, 
& 
\gTakeNotes
\\
Umgang mit Schaufeltrage, 
& 
\gTakeNotes
\\
Umgang mit Vakuummatratze, 
& 
\gTakeNotes
\\
Sandwich-Technik)
& 
\gTakeNotes
\\
- Extremitätentrauma 
& 
--- ↓ --- 
\\
(Motorik-Durchblutung-Sensibilitätskontrolle, 
& 
\gTakeNotes
\\
Stiefelgriff,
& 
\gTakeNotes
\\
Ruhigstellung, 
& 
\gTakeNotes
\\
Schienung des Armes, 
& 
\gTakeNotes
\\
Schienung des Beines, 
& 
\gTakeNotes
\\
pneumatische Schiene, 
& 
\gTakeNotes
\\
Vakuumschiene, 
& 
\gTakeNotes
\\
Extensionsschiene)
& 
\gTakeNotes
\\
- Verbandlehre
& 
\gTakeNotes
\\
- Notfälle im Säuglings- und Kleinkindalter 
&
--- ↓ --- 
\\
(Kontrolle der Lebensfunktionen und lebensrettende Sofortmaßnahmen)
& 
\gTakeNotes
\\
- An- und Auskleiden, Körperpflege und Hygiene, Nahrungs- und Flüssigkeitsaufnahme, Harn- und Stuhlentleerung,
  Erbrechen
& 
\gTakeNotes
\\
- Ergonomische und schonende Arbeitsweise (Richtiges Heben und Tragen)
& 
\gTakeNotes
\\
- Handhabung der in Einsatzfahrzeugen zu verwendenden Geräte  (insbesondere
Krankentrage, Tragsessel,  Sauerstoffgeräte, Kommunikationseinrichtungen) sowie
die  Handhabung von Rollstühlen und Gehhilfen

\\\bottomrule

\end{xtabular}